\chapter{Executive Overview}

Because of the aviation industry´s move towards net-zero global emissions in the coming decades, optimisation of aircraft configurations and performance improvements have been the focus of aircraft design in recent years. These improvements can be achieved through implementing more lightweight structures of lower stiffness, but also through certain aerodynamic features, such as increasing the wing aspect ratio, the use of winglets, or folding wing-tips. As a consequence, new challenges arise in the form of aeroservoelastic interactions between airflow, structures, rigid-body motions, and control systems. 

The aforementioned phenomena can be static or dynamic, some of the former being torsional divergence and control reversal, and some of the latter being flutter. Such effects can occur on any part of the airplane, such as the main wing, the empennage or the control surfaces. In general, they can lead to catastrophic failure and have to be avoided within the operational flight envelope. This is often achieved with the use of aeroelastic tailoring; a passive approach where the structural elements (such as composite panels) are stiffened in a specific direction, or mass balances which are attached to the wing to modify the inertial properties. While these passive approaches do mitigate the problems to a certain extent, they optimize the airplane for a specific flight condition, and often come at the cost of additional weight.

An alternative, emerging approach, is to design active control systems which prevent the onset of possibly catastrophic oscillations (active flutter suppression systems), actively alleviate loads (manoeuvre and gust load alleviation systems) and improve passenger ride comfort \cite{livne}. These systems rely on robust and adaptive control algorithms, the feedback and feedforward information from sensors, as well as control surfaces powered by high-bandwidth actuators that physically apply the corrective control inputs. The benefit of these active approaches is not only higher efficiency in terms of better performance, handling characteristics, and increase in safety, but also in terms of expanding the operational flight envelope.

Nevertheless, despite these advantages, the testing and certification of these flight control systems remains a challenge. Accurate computational simulation is difficult due to the inherent nonlinearity of high-aspect-ratio wings. On the other hand, wind tunnel tests cannot accurately model the real behaviour of such aircraft both due to scaling and because the model is constrained in the wind tunnel. This is important as rigid-body motions are a crucial component for the excitation of phenomena such as body freedom flutter. Lastly, wind tunnel testing of flexible aircraft models is not only extremely expensive, but it is sometimes also prohibited in specific situations due to the risk of damaging the wind tunnel. For instance, in case of testing the post-flutter behaviour \cite{Castro2026}. The remaining approach is flight testing, an approach that provides realistic data, but notoriously precarious for its cost and pilot safety. 

Therefore, the aim of the High-speed Unmanned General-purpose Observatory (HUGO), is to design a high-aspect-ratio turbine-powered flying test-bed which is low-cost, low-weight, and versatile. It must assist researchers in testing aeroservoelastic control systems, by having the ability to carry multiple types of scientific payloads. Consequently, the idea is to create a versatile design where the test-bed itself can be modified to be able to carry the payload. This will involve the design and analysis of modular structural and aerodynamic features, the integration of the provided VULCAN engine for reaching near-transonic speeds, performance analysis, ground operation analysis, manufacturing procedure and an analysis of environmental impact.

Similar scaled aircraft for testing aeroelastic effects and control laws have been developed before. For example, in 2005, Lockheed Martin's Skunkworks began the body freedom flutter research program, leading to the development of the X-56A, with the goal of investigating active flutter suppression, and gust load alleviation of high aspect ratio wings \footnote{\href{https://www.lockheedmartin.com/en-us/products/X-56A.html}{Lockheed Martin X-56A}(accessed on 29.04.2026))}. The NLR has also developed the scaled flight demonstrator (SFD); a scale model of the A320, aiming to supplement numerical simulations, wind tunnels, and other experimental testing means, by investigating aircraft dynamic behaviour, and then pivoting to distributed electric propulsion \footnote{\label{ref:NLR SFD},\href{https://www.nlr.org/newsroom/nieuws/maiden-flight-of-clean-sky-2s-scaled-flight-demonstrator/}{NLR Scaled Flight Demonstrator Maiden Flight}(accessed on 29.04.2026))}. Although similar to HUGO, the X-56A, and SFD are both highly specific in their implementation. Although the X-56A also includes a modular wing, it's main purpose is not to serve as a primary test bed for versatile experimentation. Additionally, neither aircraft flies above 150kts, staying well within the incompressible flow regime \footref{ref:NLR SFD} \footnote{\href{https://www.militaryfactory.com/aircraft/detail.php?aircraft_id=1115}{Military Factory Lockheed Martin X-56A}(accessed on 29.04.2026))}. HUGO provides a significant advantage as it goes faster (in the transonic regime), and would allow for testing a much wider range of experiments, aiming to replace wind tunnel and flight testing within it's niche.

Based on the user requirements and the existing configurations, the project objectives are defined. Financial objectives relate to the maximum manufacturing cost per configuration of $45000\texteuro{}$ and the maximum operating cost of $2000\texteuro{}$. Quality objectives ensure that the scientific measurements can be collected with satisfactory standards. Performance objectives relate to maximum speed, endurance and payload capacity. Business objectives specify the required profit and market share margins. Finally, compliance objectives relate to compliance with airworthiness requirements.

Using the objectives, the market is analysed to obtain a clear overview of the possible customers of such flying laboratories, as well as strengths, weaknesses, opportunities, and threats (SWOT).  Universities and research institutes, as well as aerospace companies and defence agencies have been identified as potentially interesting market segments, with the first two being identified as the most relevant for this project. The most important market opportunities identified for this project are the overcoming of the limitations of wind tunnels, and the most severe market threats relate to flight certification of the flying laboratory.

Furthermore, the functional analysis has been performed, including the mission profile definition, functional flow diagram and functional breakdown structure. The mission has been subdivided into the launch, flight test and recovery phases. Based on this, the requirement discovery tree has been created and the final list of system requirements has been compiled.

Discarding the unfeasible configurations among the numerous options \autoref{tab:subsystem_screening}, followed by the selection of the most feasible options with respect to project objectives yielded the following four configurations:

\begin{enumerate}
    \item Low wing, T-tail, retractable gear: the traditional option, mimicking standard commercial aircraft configurations
    \item Low wing, H-tail, retractable gear: offers additional possibilities for engine placement and tail-wake interference
    \item Low wing, inverted V-tail, retractable gear: offers additional possibilities for test specimen mounting, and cleaner aerodynamic flow over the top of the aircraft
    \item Flying wing with buried engine(s): low parasitic drag, simulation of coupled aeroelastic phenomena such as the Body Freedom Flutter
\end{enumerate}

In the baseline sustainability assessment, the aspects of construction, operation, end-of-life solutions, lifecycle assessment and environmental benefits have been considered. The conclusions were that manufacturing and material choice, the NOX emissions due to SAF use, recyclability, streamlined lifecycle assessment, and potential environmental benefits should all be considered during the design process.





1. Project Objectives
    - Financial
    - Quality
    - Performance
    - Business
    - Compliance

2. Market Analysis
    - Market Segmentation
    - SWOT analysis
3. Functional Analysis
    - Mission Profile
    - Functional Breakdown Structure
    - Functional Flow Diagram

4. Requirements and Constraints
    - Requirement Discovery Tree
    - Final list of system requirements

5. Resource Allocation
    - identifying the available budget
    - Centre of gravity location
    - Communication Link
    - Mass
    - Aircraft Dimensions
    - Fuel Budget
    - Power Available
    - Human Factor
    - Strategic Resource Planning
    - User Tracking

6. Preliminary Configuration Selection
    - Design Option tree
    - Generation of preliminary design concepts
    - Short Listing Configurations

7. Sustainability analysis
    - Construction
    - Operation
    - EOL solutions
    - LCA
    - Environmental benefits

The profitability and return analysis determined the size of the markets for aerospace component testing, the total cost of goods sold and the profitability of HUGO, resulting in a total attainable market share in the Netherlands for a target year. The total addressable global market was found to be \$6.94B in 2026, and \$1B in Europe. Data for the serviceable addressable market in the Netherlands was not available, so it was approximated using a bottom-up approach, resulting in \texteuro{71M}, and in a final serviceable obtainable market of \texteuro{14.2M}. As seen in \autoref{Overview-Cost-Table}, the production and direct operational costs; estimated as a an expense per flight day or per 300 flight days, are listed. These were used to calculate the cost of goods sold (COGS) per flight day of \texteuro{1142} at a price of \texteuro{2000}. For a total of 100 flight days per year, this represents a net annual profit of \texteuro{63,664} per HUGO. 


\begin{table}[H] \label{Overview-Cost-Table}
    \centering
    \caption*{Table 9.1: Overview of UAV costs and their contribution to flight day pricing}
    
    \begin{tabular}{|l|l|l|c|r|}
        \hline
        \textbf{Type} & \textbf{Subcategory} & \textbf{Description} & \textbf{Pricing Role} & \textbf{Cost (€)} \\ \hline \hline
        
        \multicolumn{5}{|c|}{\textbf{Direct costs, reference flight day}} \\ \hline
        Operations & Turbine fuel & 11.5 L of SAF, 4 test flights & Direct & 120 \\ \hline
        Operations & Electricity & Batteries, computers, systems & Direct & 80 \\ \hline
        Operations & Transport & Vehicle + logistics & Direct & 250 \\ \hline
        Operations & Mission & UAV pilot/operator & Direct & 400 \\ \hline
        
        \multicolumn{5}{|c|}{\textbf{Fractional costs, full return in 300 flight days}} \\ \hline
        Certification & RDW + ILT & Registration and flight permits & Fraction & 5\,000 \\ \hline
        
        Production & Materials & UAV structural and system components & Fraction & 28\,000 \\ \hline
        Production & Labour & 300 h at 50 €/h & Fraction & 15\,000 \\ \hline
        Production & Facility & 10 days at 200 €/day & Fraction & 2\,000 \\ \hline
        
        Operations & Insurance & 500 € annually, 15 years total & Fraction & 7\,500 \\ \hline
        
        Maintenance & Engine & 250 € per 50 h, 3000 h total & Fraction & 15\,000 \\ \hline
        Maintenance & Airframe & 500 € per 100 h, 3000 h total & Fraction & 15\,000 \\ \hline
        
    \end{tabular}
\end{table}


9. Technical Risk Assessment 
    - Initial Risk Assessment
    The first part of this chapter is basically risk assessment done the ADSEE way.
    
    Development and operation risks were identified and assessed in terms of likelihood and severity. 
    The development risks in the red zone were: discrepancies in component size/weight specified by the manufacturer vs actual values, system model errors (inadequate V\&V).
    From operation risks, the biggest concerns are related to mechanical failure of components, exceeding flight limit area, and components falling off during flight.
    
    - Contingency management
    This section is focused on lowering the likelihood and severity scores from Initial Risk Assessment.
    Lowering the likelihood was achieved by identifying preventive measures corresponding to all development risks and some operation risks.
    Lowering the severity was achieved by identifying mitigation strategies for the operational risks.
    The new risk matrix features no risks in the red zone.
    
    - Preliminary Specific Operation Risk Assessment (SORA)
    SORA is a 10 step classification methodology used for determining the operation risk of a UAV.

    For our analysis, we don't have a location yet, so it is assumed that HUGO will fly only in a controlled airspace, with very low population around. For this (idealised) case, it falls under SAIL II, but if there are any settlements whatsoever, it jumps up to SAIL V or SAIL VI.

    Based on the SAIL value, we have specific containment requirements and operational safety objectives (OSOs) that need to be complied with.

\textcolor{red}{+recommendations}



\textcolor{red}{+conclusion}

